% ─────────────────────────────────────────────────────────────────────────────
% Study 2: Online Event-Aware Residual Forecast Correction
% All numbers read from outputs/study2_final_rich_head/ and
% outputs/study2_rich_head/.  None are invented.
% ─────────────────────────────────────────────────────────────────────────────
\section{ Online Event-Aware Residual Forecast Correction}
\label{sec:study2}

\subsection{Motivation}
\label{subsec:study2_motivation}

Section~\ref{sec:study1_cf_audit} established that AI-READI does not
support a causal meal counterfactual: both the glucose-defined and
wearable-defined scenario arms fail sign agreement with matched
natural-experiment effects.
However, the frozen base forecaster is well-calibrated (one-step 80\%
prediction-interval coverage: 83.8\% train, 83.7\% validation, 83.5\% test;
Fig.~\ref{fig:study2_f1}), and its one-step residual carries a real surprise
signal that co-varies with glucose excursions.

Section ~\ref{sec:study1_cf_audit}  asks a narrower question:
\emph{can a lightweight correction head trained on streaming model-surprise
signals reduce forecast error at triggered windows, using only quantities
available at deployment time?}
The head requires no meal annotations, no scenario-channel supervision,
and no future covariates.

\subsection{Method}
\label{subsec:study2_method}

\paragraph{Frozen base.}
All Study~2 experiments use the \texttt{mamba\_stateful\_5epoch} checkpoint
(6.8\,MB; Mamba-based SSM, participant-disjoint splits).
The checkpoint is never updated.
A 5-minute rolling forecast cache is generated with
$\mathtt{anchor\_stride\_steps}=1$, producing 38.5\,M forecast rows across
train, validation, and test.

\paragraph{One-step surprise.}
At anchor $t$ the standardised surprise is
\begin{equation}
  z_t = \frac{r_t}{w_t + \varepsilon},
  \quad
  r_t = G_t - \hat{q}_{50,t-1}^{(1)},
  \quad
  w_t = \hat{q}_{90,t-1}^{(1)} - \hat{q}_{10,t-1}^{(1)},
  \label{eq:s2_zt}
\end{equation}
where $\hat{q}_{p,t-1}^{(1)}$ is the $p$-th quantile of the one-step
forecast issued at $t-1$ and $\varepsilon=10^{-3}$.
This quantity is causal: it uses only the immediately preceding forecast
and the glucose reading that just arrived.

\paragraph{Trigger rule.}
A rolling-3-anchor rule fires:
\begin{equation}
  \mathtt{trigger\_up}_t = 1
  \;\iff\;
  z_{t-2},\,z_{t-1},\,z_t > \tau_{\uparrow},
  \label{eq:s2_trigger}
\end{equation}
with a symmetric $\mathtt{trigger\_down}$ rule.
Both thresholds were locked on the validation split.

\paragraph{Event-enriched proxy.}
Because AI-READI contains no confirmed meal timestamps, an approximate
non-oracle proxy is defined as:
\begin{equation}
  \mathtt{event\_enriched} =
  \bigl(\mathtt{trigger\_up} \wedge G_t \geq 130\text{ mg/dL}\bigr)
  \;\vee\;
  \bigl(\mathtt{trigger\_down} \wedge G_t \leq 85\text{ mg/dL}\bigr).
  \label{eq:s2_proxy}
\end{equation}
This proxy identifies triggers that coincide with likely hyperglycaemia or
hypoglycaemia; 36.5\% of triggered non-insulin test anchors qualify.

\paragraph{Evaluation protocol.}
All metrics are restricted to non-insulin participants (validation: $n=221$;
test: $n=204$) and triggered anchors only.
Confidence intervals use participant-clustered bootstrap resampling
($B=1000$, $\alpha=0.05$).
Paired difference CIs compare methods on the same anchors.
The test split was held out during all threshold and hyperparameter selection.

\subsection{Trigger Calibration and Threshold Selection}
\label{subsec:study2_trigger}

The threshold grid $\tau \in \{0.2, 0.3, \ldots, 0.8\}$ was evaluated on
the validation split.
The selection criterion was the highest $\tau$ achieving a trigger-up rate
near 5\% while retaining all 221 non-insulin validation participants
(sufficient coverage for participant-clustered CIs).
Only $\tau=0.2$ satisfied this criterion:
trigger-up rate 5.06\%, trigger-down rate 4.38\%, all 221 participants
triggered at least once.
At $\tau \geq 0.3$ participant coverage drops below 100\% within the tested
grid (Fig.~\ref{fig:study2_f1}).

The locked threshold $\tau_{\uparrow} = \tau_{\downarrow} = 0.2$ produces
the following trigger counts (non-insulin participants):
\begin{itemize}
  \item Train: 190{,}893 triggered anchors, 1{,}043 participants, trigger rate 9.19\%.
  \item Validation: 44{,}417 anchors, 221 participants, trigger rate 9.42\%.
  \item Test: 39{,}698 anchors, 204 participants, trigger rate 9.26\%.
\end{itemize}
Rates are stable across splits, indicating the trigger is not overfit.

\subsection{Rich Causal Residual-Correction Head}
\label{subsec:study2_richhead}

\paragraph{Feature set.}
Each triggered anchor $t$ contributes the following causal trailing
features to the correction head.
All features are computed from quantities available strictly before or
at time $t$; no future glucose, future label, or future covariate is used.

\begin{itemize}
  \item \textbf{Residual history}: $r_t,\,r_{t-1},\,r_{t-2},\,r_{t-3}$
        (lagged one-step residuals, mg/dL).
  \item \textbf{Standardised surprise history}: $z_t,\,z_{t-1},\,z_{t-2}$.
  \item \textbf{Glucose slopes}:
        slope$_{15}=(G_t - G_{t-3})/15$,
        slope$_{30}=(G_t - G_{t-6})/30$,
        slope$_{60}=(G_t - G_{t-12})/60$ (mg/dL/min).
  \item \textbf{Causal curvature}: $(\text{slope}_{15} - \text{slope}_{30})/15$ (mg/dL/min$^2$).
  \item \textbf{Interval width}: $w_t = \hat{q}_{90} - \hat{q}_{10}$ (model uncertainty, mg/dL).
  \item \textbf{Current glucose}: $G_t$ (mg/dL).
  \item \textbf{Trigger-phase features}: time-since-trigger onset (anchors),
        cumulative residual since trigger onset (mg/dL).
  \item \textbf{Circadian proxy}: $\sin(2\pi h/24)$ and $\cos(2\pi h/24)$,
        where $h = \mathtt{hours\_since\_start} \bmod 24$.
  \item \textbf{Direction flag}: $\mathtt{trigger\_up} \in \{0,1\}$.
\end{itemize}
Two features were excluded as zero-variance after inspection:
the column \texttt{steps\_since\_start} is a sequential anchor counter
(not a pedometer), so its 15- and 30-minute differences are constant.
Heart-rate change features are absent from the cache parquet.
The final feature set contains \textbf{19 features}.

\paragraph{Architecture and training.}
Separate correction heads are trained for upward and downward triggers.
Each head consists of 12 independent Ridge regression models
($\alpha=1.0$, with intercept, features standardised),
one per forecast horizon $h \in \{1,\ldots,12\}$.
The target for horizon $h$ is $G_{t+h} - \hat{q}_{50,t}^{(h)}$, the
correction the head must add to the frozen-base median.
Both up- and down-heads are trained on the non-insulin train split
(190{,}893 triggered anchors); all 1{,}043 train participants contribute
triggers in both directions.

\paragraph{Head selection.}
Candidate heads and baselines were evaluated on the validation split only.
The two top-ranked methods -- Ridge full (split up/down) and HGBT full
(combined) -- were statistically tied (validation paired difference CI:
$-0.0009$ mg/dL, 95\% CI $[-0.040,\;+0.029]$, CI includes zero).
Following a pre-specified tie-break rule of preferring the simpler model,
\textbf{Ridge full (split up/down)} was selected as the final head.
Test metrics were examined only after this selection was fixed.

\subsection{Baseline Ladder}
\label{subsec:study2_ladder}

Table~\ref{tab:study2_ladder_rich} and Fig.~\ref{fig:study2_f2_rich} present
MAE with participant-clustered 95\% CIs for all methods on non-insulin
triggered anchors.
Anchor counts are reported as distinct $(participant\_id,\;segment\_id,\;anchor\_ds)$
tuples; the frozen-base anchor is never counted more than once per participant
and anchor index.

\paragraph{Key findings.}
\begin{itemize}
  \item \textbf{Current shift is worse than the frozen base} ($+0.217$~mg/dL
        on test).
        Applying the full residual $r_t$ as a constant offset overcorrects:
        the base model already mean-reverts toward the historical glucose
        distribution over longer horizons, and a flat shift counteracts this.
  \item \textbf{Decaying shift matches the frozen base}: test MAE
        $11.381$ vs $11.383$~mg/dL (difference negligible).
  \item \textbf{Linear $z_t$ beats decaying shift} ($-0.025$~mg/dL).
        Standardising by the model's own uncertainty width $w_t$ provides
        a better correction scale than a fixed decay.
  \item \textbf{The rich head (Ridge split) beats the best simple baseline}
        (linear $z_t$ + slope + curvature) by $0.122$~mg/dL (test),
        95\% CI $[-0.163,\;-0.091]$ (CI excludes zero).
  \item \textbf{Ridge and HGBT are statistically tied}: HGBT combined minus
        Ridge split $= -0.005$~mg/dL, 95\% CI $[-0.029,\;+0.023]$
        (CI includes zero).
\end{itemize}

\paragraph{Ablation.}
Fig.~\ref{fig:study2_f6_ablation} shows validation MAE for Ridge trained on
six feature groups.
Residual history ($r_t$--$r_{t-3}$) is the strongest single group
(val MAE $11.340$, test $11.302$), outperforming $z_t$ alone (val $11.403$).
The full 19-feature set achieves val MAE $11.290$, lower than any single group
in isolation, confirming complementary information across feature families.

\subsection{Event-Enriched Proxy versus Unmatched Control}
\label{subsec:study2_event}

\emph{This comparison is descriptive, not a matched causal contrast.
Event-enriched labels are glucose-level proxies because AI-READI does not
contain confirmed meal timestamps.
The control group is not distribution-matched on glucose level, slope,
$z_t$, or interval width.
The structural confound arises from the proxy definition itself:
Equation~\eqref{eq:s2_proxy} selects on glucose level, so event and
non-event anchors differ by construction.}

Table~\ref{tab:study2_event_rich} and Fig.~\ref{fig:study2_f3_rich} compare
correction gain (base MAE minus final-head MAE) for the event-enriched and
unmatched-control groups on the test split.

\begin{itemize}
  \item \textbf{Event-enriched} ($n=14{,}505$ anchors):
        base MAE $14.355$, head MAE $13.941$,
        gain $+0.415$~mg/dL,
        95\% CI $[+0.216,\;+0.603]$ (CI excludes zero).
  \item \textbf{Unmatched control} ($n=25{,}193$ anchors):
        base MAE $9.671$, head MAE $9.635$,
        gain $+0.036$~mg/dL,
        95\% CI $[+0.003,\;+0.070]$ (CI excludes zero).
\end{itemize}

The event group shows larger absolute gain, but this is expected:
the event definition selects anchors with high glucose, which are
inherently harder to forecast (higher base MAE) and where large
one-step residuals are more common.
The control group gain, while small, also excludes zero, indicating
that the correction head is not harmful outside event windows.
Neither comparison should be interpreted as evidence that the head
detects or responds specifically to physiological events.

\subsection{Deployment Behavior and Abstention}
\label{subsec:study2_deploy}

Fig.~\ref{fig:study2_f5} and Table~\ref{tab:study2_deploy} summarise
deployment characteristics across splits.

\begin{itemize}
  \item \textbf{Trigger rate}: 9.19\% (train), 9.42\% (validation),
        9.26\% (test) of non-insulin anchors.
        Stability across splits confirms the trigger is not overfit.
  \item \textbf{Episode duration}: median 5 minutes, P90 15 minutes
        in all splits.
        With $\tau=0.2$ the rolling-3 rule rarely sustains for more than
        one or two consecutive anchors.
  \item \textbf{Abstention rate}: $< 0.09\%$ in all splits, arising
        exclusively from stream boundaries (first anchor of each segment)
        and rare missing prediction-interval rows ($< 0.001\%$).
\end{itemize}

\subsection{Interpretation and Limitations}
\label{subsec:study2_interp}

\paragraph{What the head achieves.}
The rich correction head reliably reduces MAE at triggered windows.
The improvement over the frozen base ($0.175$~mg/dL test, 95\% CI
$[-0.259,\;-0.094]$) and over the best non-trivial no-learning baseline
(decaying shift: $0.172$~mg/dL test, 95\% CI $[-0.238,\;-0.109]$)
are both statistically clear.
The improvement over the best simple linear baseline ($0.122$~mg/dL,
CI $[-0.163,\;-0.091]$) is also statistically clear, showing that the
richer causal feature set adds genuine value beyond a single $z_t$ term.
The ablation confirms that residual history ($r_t$--$r_{t-3}$) is the
most informative feature group, followed by trigger-phase features
and the combined set of slopes and curvature.

\paragraph{What the head does not achieve.}
The gap between Ridge and HGBT is not statistically distinguishable
($-0.005$~mg/dL, CI $[-0.029,\;+0.023]$), indicating that the linear
form of Ridge already captures the dominant structure in this feature space.
The event-vs-control comparison is descriptive and structurally confounded;
it should not be read as evidence that the head specifically targets
physiological events.
AI-READI lacks confirmed meal or exercise timestamps, so no oracle event
evaluation is possible.

\paragraph{Per-horizon profile.}
Fig.~\ref{fig:study2_f7_horizon} shows per-horizon MAE for the final
head and selected baselines.
The correction provides the largest absolute benefit at short horizons
(0--30 min): the rich head reduces 0--30 min MAE from
$8.112$ (base) to $7.844$~mg/dL.
The benefit narrows at long horizons: 30--60 min MAE is $14.573$~mg/dL
(head) versus $14.654$~mg/dL (base).
This pattern is expected because the one-step residual is most informative
about the immediate future, and its signal decays with horizon.

\paragraph{Scope.}
The correction head is appropriate only when triggered
($\approx 9.3\%$ of deployment anchors).
At non-triggered anchors the frozen-base median should be used.
The head should be retrained if the base checkpoint is updated or if
the target glucose population changes substantially.

\subsection{Link to Interpretability}
\label{subsec:study2_link_interp}

The correction head's 19-feature linear Ridge structure remains interpretable
by design: each coefficient $b_{h,j}$ quantifies how much feature $j$ shifts
the horizon-$h$ prediction.
The dedicated interpretability section now focuses on the frozen base model's
two reliance views: output-level LOMO attribution and rigorous hidden-state
construction through $h_t$.
This separates base-model reliance from the Study~2 correction-head audit.

% ── Figures ──────────────────────────────────────────────────────────────────

\begin{figure}[H]
  \centering
  \includegraphics[width=\linewidth]{generated/study2_f1_trigger_calibration.pdf}
  \caption{%
    \textbf{S2-F1: Trigger calibration.}
    (a) Empirical one-step 80\% PI coverage by split;
    the dashed line marks the nominal 80\% target.
    (b) Trigger-up rate as a function of threshold $\tau$;
    the locked threshold $\tau=0.2$ is marked.
    (c) Trigger-up and trigger-down rates at $\tau=0.2$ by split.
    Coverage exceeds 83\% in all splits; $\tau=0.2$ is the only candidate
    achieving near-5\% trigger rate with 100\% participant coverage.%
  }
  \label{fig:study2_f1}
\end{figure}

\begin{figure}[H]
  \centering
  \includegraphics[width=\linewidth]{generated/study2_f2_baseline_ladder_rich.pdf}
  \caption{%
    \textbf{S2-F2: Full correction-head baseline ladder.}
    MAE (mg/dL) with participant-clustered 95\% CIs for all ten correction
    methods on non-insulin triggered windows.
    Left: validation (44{,}417 anchors, $n=221$);
    right: test (39{,}698 anchors, $n=204$).
    Dashed vertical line: frozen-base MAE.
    Shaded band: gain region of the selected final head (Ridge split up/down).
    Anchor counts are distinct $(participant, segment, anchor)$ tuples.%
  }
  \label{fig:study2_f2_rich}
\end{figure}

\begin{figure}[H]
  \centering
  \includegraphics[width=\linewidth]{generated/study2_f8_paired_cis.pdf}
  \caption{%
    \textbf{S2-F8: Paired-CI forest plot (test split).}
    Each bar shows the gain of the richer method over the comparison method,
    with participant-clustered 95\% CI.
    Blue: CI excludes zero (statistically distinguishable).
    Grey: CI includes zero (statistically tied).
    Asterisk (*): CI excludes zero.
    All comparisons against simpler baselines favour the rich head;
    HGBT and Ridge are indistinguishable.%
  }
  \label{fig:study2_f8}
\end{figure}

\begin{figure}[H]
  \centering
  \includegraphics[width=\linewidth]{generated/study2_f3_event_control_rich.pdf}
  \caption{%
    \textbf{S2-F3: Event-enriched proxy versus unmatched control (test split).}
    Left: MAE of the frozen base and final head for each group.
    Right: correction gain (base MAE $-$ head MAE) with 95\% CI.
    \emph{Descriptive only -- not a matched causal contrast.
    Event-enriched labels are glucose-level proxies because AI-READI does
    not contain confirmed meal timestamps.
    The control group is not distribution-matched on glucose level, slope,
    $z_t$, or interval width.}%
  }
  \label{fig:study2_f3_rich}
\end{figure}

\begin{figure}[H]
  \centering
  \includegraphics[width=0.85\linewidth]{generated/study2_f6_ablation.pdf}
  \caption{%
    \textbf{S2-F6: Feature-group ablation (Ridge, validation MAE).}
    Each bar shows validation MAE for a Ridge head trained on the named
    feature subset only (19-feature full-causal set in green).
    Error bars: participant-clustered 95\% CI.
    Dashed vertical line: frozen-base validation MAE.
    The full causal set achieves the lowest validation MAE;
    residual history ($r_t$--$r_{t-3}$) is the strongest single group.%
  }
  \label{fig:study2_f6_ablation}
\end{figure}

\begin{figure}[H]
  \centering
  \includegraphics[width=0.85\linewidth]{generated/study2_f7_per_horizon.pdf}
  \caption{%
    \textbf{S2-F7: Per-horizon MAE (test split).}
    Lines show mean absolute error at each 5-minute forecast horizon
    for the frozen base, best simple baseline, and the two best rich heads.
    The correction provides the largest benefit at short horizons (0--30 min)
    and narrows at longer horizons.%
  }
  \label{fig:study2_f7_horizon}
\end{figure}

\begin{figure}[H]
  \centering
  \includegraphics[width=\linewidth]{generated/study2_f5_deployment.pdf}
  \caption{%
    \textbf{S2-F5: Deployment behavior (non-insulin participants).}
    (a) Trigger rate by split.
    (b) Episode duration: median and P90.
    (c) Abstention rate.
    Trigger rates and episode durations are stable across splits;
    abstention is negligible ($<0.09\%$ in all splits).%
  }
  \label{fig:study2_f5}
\end{figure}

% ── Tables ───────────────────────────────────────────────────────────────────

\begin{table}[H]
  \centering
  \caption{%
    \textbf{S2-T1: Final model selection (validation-based).}
    All candidate models ranked by validation MAE (lower is better).
    Anchor counts are distinct $(participant, segment, anchor)$ tuples.
    The paired difference CI (best vs.\ second-best on validation) did not
    exclude zero, so the simpler model (Ridge split up/down) was selected.%
  }
  \label{tab:study2_selection}
  \small
  \begin{tabular}{llcccc}
    \toprule
    Rank & Method & $n_{\text{anchors}}$ & Val MAE & Val 95\% CI & Selected \\
    \midrule
    1 & Ridge full (split up/down)         & 44{,}417 & 11.228 & [10.805, 11.709] & \checkmark \\
    2 & HGBT full (combined)               & 44{,}417 & 11.229 & [10.796, 11.725] & \\
    3 & HGBT full (split up/down)          & 44{,}417 & 11.236 & [10.807, 11.733] & \\
    4 & Ridge full (combined)              & 44{,}417 & 11.290 & [10.844, 11.796] & \\
    5 & Linear full causal OLS             & 44{,}417 & 11.291 & [10.844, 11.796] & \\
    6 & Linear $z_t$ + slope + curv        & 44{,}417 & 11.376 & [10.927, 11.886] & \\
    \midrule
    \multicolumn{6}{l}{\textit{Tie-break: paired CI (Ridge split $-$ HGBT combined) $= -0.0009$,}} \\
    \multicolumn{6}{l}{\textit{95\% CI $[-0.040,\;+0.029]$ (includes zero) $\Rightarrow$ prefer simpler.}} \\
    \bottomrule
  \end{tabular}
\end{table}

\begin{table}[H]
  \centering
  \caption{%
    \textbf{S2-T2: Full baseline ladder.}
    MAE (mg/dL) with participant-clustered 95\% CIs for all ten correction
    methods on non-insulin triggered windows.
    ``n'' is distinct $(participant, segment, anchor)$ tuples.
    Starred paired differences in the lower panel are CIs excluding zero.%
  }
  \label{tab:study2_ladder_rich}
  \small
  \setlength{\tabcolsep}{4.5pt}
  \begin{tabular}{llcccc}
    \toprule
    Split & Method & $n$ & $n_{\text{part.}}$ & MAE & 95\% CI \\
    \midrule
    \multirow{10}{*}{Val.}
      & Base (frozen SSM)              & 44{,}417 & 221 & 11.474 & [10.980, 12.016] \\
      & Current shift                  & 44{,}417 & 221 & 11.557 & [11.156, 12.026] \\
      & Decaying shift ($\delta=0.7$)  & 44{,}417 & 221 & 11.434 & [10.969, 11.958] \\
      & Linear $z_t$                   & 44{,}417 & 221 & 11.403 & [10.949, 11.916] \\
      & Linear $z_t$ + slope + curv    & 44{,}417 & 221 & 11.376 & [10.927, 11.886] \\
      & Linear full causal OLS         & 44{,}417 & 221 & 11.291 & [10.844, 11.796] \\
      & Ridge full (combined)          & 44{,}417 & 221 & 11.290 & [10.844, 11.796] \\
      & Ridge full (split up/down)     & 44{,}417 & 221 & 11.228 & [10.805, 11.709] \\
      & HGBT full (combined)           & 44{,}417 & 221 & 11.229 & [10.796, 11.725] \\
      & HGBT full (split up/down)      & 44{,}417 & 221 & 11.236 & [10.807, 11.733] \\
    \midrule
    \multirow{10}{*}{Test}
      & Base (frozen SSM)              & 39{,}698 & 204 & 11.383 & [10.871, 11.968] \\
      & Current shift                  & 39{,}698 & 204 & 11.599 & [11.132, 12.106] \\
      & Decaying shift ($\delta=0.7$)  & 39{,}698 & 204 & 11.381 & [10.886, 11.939] \\
      & Linear $z_t$                   & 39{,}698 & 204 & 11.355 & [10.874, 11.891] \\
      & Linear $z_t$ + slope + curv    & 39{,}698 & 204 & 11.330 & [10.850, 11.856] \\
      & Linear full causal OLS         & 39{,}698 & 204 & 11.284 & [10.812, 11.807] \\
      & Ridge full (combined)          & 39{,}698 & 204 & 11.284 & [10.812, 11.807] \\
      & Ridge full (split up/down)     & 39{,}698 & 204 & 11.208 & [10.748, 11.718] \\
      & HGBT full (combined)           & 39{,}698 & 204 & 11.204 & [10.752, 11.725] \\
      & HGBT full (split up/down)      & 39{,}698 & 204 & 11.218 & [10.764, 11.736] \\
    \midrule
    \multicolumn{6}{l}{\textit{Selected paired differences (test, participant-clustered bootstrap, $B=1000$):}} \\[2pt]
    & Final head vs base              & \multicolumn{4}{l}{$-0.174$ mg/dL, CI $[-0.259,\;-0.094]$*} \\
    & Final head vs decaying shift    & \multicolumn{4}{l}{$-0.172$ mg/dL, CI $[-0.238,\;-0.109]$*} \\
    & HGBT combined vs best simple    & \multicolumn{4}{l}{$-0.126$ mg/dL, CI $[-0.163,\;-0.091]$*} \\
    & HGBT combined vs linear OLS     & \multicolumn{4}{l}{$-0.080$ mg/dL, CI $[-0.110,\;-0.052]$*} \\
    & HGBT combined vs Ridge split    & \multicolumn{4}{l}{$-0.005$ mg/dL, CI $[-0.029,\;+0.023]$ (tied)} \\
    \bottomrule
    \multicolumn{6}{l}{\small * CI excludes zero.} \\
  \end{tabular}
\end{table}

\begin{table}[H]
  \centering
  \caption{%
    \textbf{S2-T3: Event-enriched proxy versus unmatched control (test split).}
    Correction gain (base MAE $-$ final-head MAE, mg/dL) for event-enriched
    and unmatched-control groups.
    CIs are participant-clustered bootstrap ($B=1000$).
    \emph{This comparison is descriptive, not a matched causal contrast.
    The control group is not distribution-matched on glucose level, slope,
    $z_t$, or interval width.}%
  }
  \label{tab:study2_event_rich}
  \small
  \begin{tabular}{lccccc}
    \toprule
    Group & $n$ & Base MAE & Head MAE & Gain & 95\% CI \\
    \midrule
    Event-enriched  & 14{,}505 & 14.355 & 13.941 & $+0.415$ & $[+0.216,\;+0.603]$ \\
    Unmatched control & 25{,}193 & 9.671 & 9.635  & $+0.036$ & $[+0.003,\;+0.070]$ \\
    \bottomrule
  \end{tabular}
\end{table}

\begin{table}[H]
  \centering
  \caption{%
    \textbf{S2-T4: Deployment characteristics (non-insulin participants).}
    Trigger rates and episode statistics are stable across splits.%
  }
  \label{tab:study2_deploy}
  \small
  \begin{tabular}{lcccccc}
    \toprule
    Split & $n_{\text{part.}}$ & Trigger rate & Episodes
      & Dur.\ med. & Dur.\ P90 & Abstention \\
    \midrule
    Train      & 1{,}043 & 9.19\% & 111{,}288 & 5\,min & 15\,min & 0.083\% \\
    Validation &   221   & 9.42\% &  25{,}244 & 5\,min & 15\,min & 0.078\% \\
    Test       &   204   & 9.26\% &  23{,}035 & 5\,min & 15\,min & 0.083\% \\
    \bottomrule
  \end{tabular}
\end{table}

\begin{table}[H]
  \centering
  \caption{%
    \textbf{S2-T5: Feature-group ablation (Ridge, validation MAE).}
    Each row shows validation and test MAE for a Ridge head
    trained on the named feature subset only.
    The full 19-feature set achieves the lowest validation MAE.
    All $n_{\text{anchors}}$ are the same 44{,}417 validation and
    39{,}698 test triggered tuples.%
  }
  \label{tab:study2_ablation}
  \small
  \begin{tabular}{lccccc}
    \toprule
    Feature group & $n_{\text{feat.}}$ & Val MAE & Val 95\% CI & Test MAE \\
    \midrule
    $z_t$ only                         & 1  & 11.403 & [10.949, 11.916] & 11.355 \\
    $z_t$ + slope + curv               & 3  & 11.376 & [10.927, 11.886] & 11.330 \\
    $r_t$ only                         & 1  & 11.356 & [10.908, 11.859] & 11.310 \\
    Slope + curvature                  & 4  & 11.401 & [10.942, 11.909] & 11.349 \\
    Trigger phase                      & 4  & 11.348 & [10.902, 11.853] & 11.314 \\
    Residual history ($r_t$--$r_{t-3}$)& 4  & 11.340 & [10.896, 11.844] & 11.302 \\
    Context (glucose, IW, hour)        & 4  & 11.455 & [10.988, 11.977] & 11.383 \\
    \textbf{Full causal (19 features)} & 19 & \textbf{11.290} & [10.844, 11.796] & \textbf{11.284} \\
    \bottomrule
  \end{tabular}
\end{table}
